\section{Benchmark}
\label{sec:benchmark}

\subsection{Robots and task}
\label{sec:robots}

\paragraph{Wojtek} is a small quadruped built by the authors, with a
standing height of about 0.10\,m and 12 actuated joints (three per leg).
Each leg is a closed kinematic chain: the actuated joints drive the foot
through a passive four-bar linkage, closed in MuJoCo by one
\texttt{connect} equality constraint per leg. The closure adds constraint
dynamics that serial legs do not have (the linkage can reach a toggle
configuration, which the reward penalises), which is why we pair Wojtek
with an open-chain robot.
\todo{Add mass, link lengths and a photo of the physical robot.}

\paragraph{Unitree Go1} is the open serial-chain quadruped of the MuJoCo
Playground joystick task~\citep{zakka2025playground}, with Playground's
reward, feet-only collisions and physics randomisation. Its standing trunk
height is 0.258\,m. We change four things so it runs the same benchmark:
the scene is our heightfield arena; commands are forward-walking;
termination uses Wojtek's fall rule at Go1's size (trunk below
$0.6\times0.258$\,m above the terrain under it, or gravity $z$-component in
the body frame above $-0.4$); and the critic receives a terrain height scan.

\paragraph{Task.}
Both robots track a velocity command
$(v_x, v_y, \omega_z)$ drawn per episode from
$v_x\in[0.2,0.6]$, $v_y\in[-0.2,0.2]$\,m/s and
$\omega_z\in[-0.4,0.4]$\,rad/s for Wojtek; Go1's linear ranges are scaled
by $s=2.5$. The control step is 0.02\,s and an episode on a level lasts
300 steps (6\,s). The actor is blind (no terrain perception); the
critic additionally sees a clean $5\times5$ height scan ahead of the robot
(asymmetric actor--critic), so value-based level scores can distinguish
terrains. Training applies physics randomisation~\citep{tobin2017domain}
per environment (Wojtek: terrain friction, base and link masses,
actuator gain and damping scale; Go1: Playground's randomisation).

\paragraph{Arena.}
Levels are rendered into a square heightfield around a flat spawn disc.
In MJX the heightfield data is an array leaf of the model, so each
parallel environment can be reset onto a different level without
recompilation, which makes per-environment level swapping feasible at RL
scale (details in Appendix~\ref{app:render}).

\paragraph{Scaling.}
We set $s = 2.5$, approximately the ratio of Go1's standing height to
Wojtek's ($0.258/0.10$).
Every length in the benchmark (arena, terrain features, gap depth and
period, commands, height-scan grid and success distance) is multiplied by
$s$ for Go1; slope angles are unchanged. Capability envelopes are reported
in Wojtek units so the two robots can be compared directly.

\subsection{Level space}
\label{sec:levelspace}

A level is a vector $\theta\in[0,1]^8$ plus a noise seed. Each coordinate
maps linearly to a physical range (Table~\ref{tab:levelspace}). Features are
measured radially from the edge of the spawn disc, or are seeded random
fields, so the level does not depend on the robot's heading. The rendered
height is the sum of a radial slope, radial stairs, smooth noise, and
piecewise-constant blocks, with gap trenches (0.15\,m deep, every 0.6\,m)
cut in. Each parameter has a \emph{flat} value (zero slope, zero step,
zero amplitude, zero gap, zero block height); because slope and step ranges
are asymmetric, flat is not the centre of the cube.

\begin{table}[t]
  \centering\small
  \caption{Level space (v10), Wojtek units; Go1 multiplies every length by
  $s=2.5$. A level is $\theta\in[0,1]^8$ plus a noise seed.}
  \label{tab:levelspace}
  \begin{tabular}{lll}
    \toprule
    Parameter & Range & Pilot (v8) range \\
    \midrule
    Noise amplitude & 0--0.10\,m & same \\
    Noise wavelength & 0.2--1.0\,m & same \\
    Radial slope & $-20^\circ$ to $+18^\circ$ & $\pm20^\circ$ \\
    Stair height & $-0.08$ to $+0.05$\,m & $\pm0.08$\,m \\
    Stair width & 0.15--0.60\,m & same \\
    Gap width & 0--0.06\,m & 0--0.25\,m \\
    Block height & 0--0.06\,m & 0--0.10\,m \\
    Block size & 0.15--0.60\,m & same \\
    \bottomrule
  \end{tabular}
\end{table}

\paragraph{Calibration.}
A level space is only useful for comparing curricula if uniform sampling is
neither trivially easy (nothing to curate) nor mostly unsolvable (curation
and DR both train on impossible levels). Our pilot space (v8) was the
latter: the best pilot policies solved about 8\% of uniform levels, and
about 92\% of uniform levels were unsolved by every method. We recalibrate
with the pilot policies themselves: for each feature we measure the best
v8 policies' capability envelope, the largest feature value still passed
with 50\% success when that feature alone is active
(\S\ref{sec:protocol}), and cap the uphill slope, step-up height, gap width
and block height at that envelope. Estimated from the subset of v8
evaluation levels that fall inside the new box, v8 policies solve about
39\% of uniform v10 levels. Because features compound, uniform levels still
range from trivial to unsolved. We regard this procedure, sizing the space
to the capability envelopes of a reference policy, as part of the protocol:
it is cheap, uses only the evaluation tooling, and transfers to any
parametric level space.
\todo{Report the v10 uniform success of DR as a check that the calibration
target was met, once results land.}
